\chapter{Conception et architecture du systeme}
\label{chap:conception}

\section{Architecture globale}

La plateforme est une application web monolithique modulaire, organisee en couches, et
deployee sous forme de conteneurs Docker orchestres par Docker Compose. Le
\cref{fig:deployment} presente l'architecture de deploiement retenue : un conteneur
applicatif Spring Boot, une base de donnees relationnelle PostgreSQL, et un moteur
d'inference local (Ollama) charge du modele de langage utilise par l'assistant RAG.
Un conteneur ephemere (\emph{ollama-init}) telecharge automatiquement le modele au
premier demarrage, garantissant un environnement reproductible sans dependance a un
service LLM externe payant.

\begin{figure}[H]
\centering
\begin{tikzpicture}[
  box/.style={draw,rounded corners,minimum width=3.6cm,minimum height=1.6cm,align=center,font=\small,fill=dxcblue!6},
  cyl/.style={cylinder,shape border rotate=90,draw,aspect=0.25,minimum height=1.6cm,minimum width=2.6cm,align=center,font=\small,fill=dxcgray!10},
  arr/.style={-{Latex[length=2mm]},thick}
]
\node[box] (browser) at (0,3.2) {Navigateur\\ (HTMX / Chart.js\\ ou client Angular)};
\node[draw,dashed,rounded corners,minimum width=11.6cm,minimum height=4.4cm] (host) at (5.8,0) {};
\node[font=\bfseries\small] at (5.8,2.05) {Hote Docker (docker-compose.yml)};

\node[box] (backend) at (2.4,0.3) {\textbf{app-backend}\\ Spring Boot 3 / Java 17\\ port 8080};
\node[cyl] (db) at (6.3,0.3) {\textbf{postgres\_db}\\ PostgreSQL 15\\ port 5432};
\node[box] (ollama) at (10.2,0.3) {\textbf{ollama\_engine}\\ moteur LLM (GPU)\\ port 11434};
\node[box,minimum height=1cm,fill=dxcgray!6] (init) at (10.2,-1.9) {\scriptsize \textbf{ollama\_downloader}\\ \scriptsize pull mistral:7b (init)};

\node[cyl,minimum width=2.2cm] (vol1) at (2.4,-1.9) {\scriptsize volume\\ \scriptsize uploads};

\draw[arr] (browser.south) -- (backend.north);
\draw[arr] (backend.east) -- (db.west);
\draw[arr] (backend.north) to[out=55,in=125] (ollama.north);
\draw[arr] (init.north) -- (ollama.south);
\draw[arr] (backend.south) -- (vol1.north);
\end{tikzpicture}
\caption{Diagramme de deploiement (conteneurs Docker Compose)}
\label{fig:deployment}
\end{figure}

Au sein du conteneur applicatif, le code est organise selon une architecture en
couches classique (\cref{fig:layers}) : les \emph{controllers} exposent les points
d'entree HTTP (vues Thymeleaf/HTMX pour l'interface, endpoints REST pour le flux SSE
et l'API), les \emph{services} portent la logique metier (calcul du risque,
agregation des metriques, orchestration des jobs Spring Batch), les
\emph{repositories} (Spring Data JPA) assurent la persistance, et les
\emph{entities} modelisent le domaine. Des \emph{DTO} et des \emph{mappers}
(MapStruct) isolent le modele de domaine des formats d'entree/sortie (XML, JSON,
formulaires HTML). La communication entre composants faiblement couples (fin
d'ingestion, soumission d'une configuration, fin de recalcul) transite par un bus
d'evenements applicatif (\texttt{ApplicationEventPublisher}), ce qui decouple le
declencheur d'un recalcul de son execution effective.

\begin{figure}[H]
\centering
\begin{tikzpicture}[
  layer/.style={draw,minimum width=11cm,text width=10.6cm,minimum height=1cm,align=center,font=\small},
  node distance=6mm
]
\node[layer,fill=dxcblue!25] (l1) {Presentation -- Controllers (MVC Thymeleaf/HTMX, REST, SSE) \\ {\scriptsize\itshape via DTO et Mapper MapStruct}};
\node[layer,fill=dxcblue!18,below=of l1] (l2) {Metier -- Services (RiskCalculationService, DashboardMetricsService, JobService, AiAgentService, ...) \\ {\scriptsize\itshape couples faiblement via ApplicationEventPublisher}};
\node[layer,fill=dxcblue!12,below=of l2] (l3) {Traitement par lots -- Spring Batch (Readers StAX, Processors, JpaItemWriters)};
\node[layer,fill=dxcblue!8,below=of l3] (l4) {Acces aux donnees -- Repositories Spring Data JPA};
\node[layer,fill=dxcgray!12,below=of l4] (l5) {Persistance -- PostgreSQL};
\draw[-{Latex[length=1.8mm]}] (l1.south) -- (l2.north);
\draw[-{Latex[length=1.8mm]}] (l2.south) -- (l3.north);
\draw[-{Latex[length=1.8mm]}] (l2.south west) to[out=200,in=160] (l4.west);
\draw[-{Latex[length=1.8mm]}] (l3.south) -- (l4.north);
\draw[-{Latex[length=1.8mm]}] (l4.south) -- (l5.north);
\end{tikzpicture}
\caption{Architecture logicielle en couches}
\label{fig:layers}
\end{figure}

\section{Modele de donnees -- diagramme de classes}

Le \cref{fig:class-diagram} presente le coeur du modele de domaine. Un \texttt{Host}
(hote scanne par Qualys) porte une collection de \texttt{VulnInfo} (relation de
composition : la suppression d'un hote entraine celle de ses vulnerabilites,
\texttt{CascadeType.ALL}). Un \texttt{CmdbCi} (element de configuration CMDB) n'est
pas relie a \texttt{Host} par une cle etrangere, mais par une \emph{reconciliation
applicative} sur l'adresse IP, effectuee au moment du calcul de risque -- ce choix,
discute en \cref{sec:reconciliation}, permet d'ingerer independamment les deux
sources sans contrainte d'integrite referentielle rigide entre deux systemes
externes de rythmes de mise a jour differents. \texttt{RiskConfig} est une
configuration singleton portant les parametres fixes de la formule (poids de
severite, facteurs de modulation, seuils SLA) ; \texttt{RiskWeightConfig} est une
table generique \texttt{(categorie, cle) -> poids} qui porte les poids dont
l'ensemble des cles est variable (roles techniques, environnements, classes de
risque pour le MTTR). \texttt{RiskSnapshot} historise, a chaque recalcul, la
repartition du parc par classe et l'ARS moyen, alimentant le graphique d'historique
du tableau de bord.

\begin{figure}[H]
\centering
\begin{adjustbox}{width=\textwidth}
\begin{tikzpicture}[
  cls/.style={draw,rectangle split,rectangle split parts=3,rectangle split part align={left,left,left},
              font=\scriptsize,text width=4.1cm,rectangle split part fill={dxcblue!20,white,white}},
  rel/.style={draw,-},
  dep/.style={draw,dashed,-{Latex[length=2.2mm,width=1.6mm]}},
  note/.style={draw,dxcgray,fill=yellow!8,align=left,font=\scriptsize,text width=3.4cm,inner sep=4pt},
  noteline/.style={draw,dxcgray,dashed},
]
\node[cls] (host) at (0,0) {
  \textbf{Host}
  \nodepart{second} - id: Long \\ - ip: String \\ - dns: String \\ - os: Os \\ - acs: Double \\ - ars: Double \\ - riskClass: RiskClass \\ - slaStatus: SlaStatus \\ - slaDueDate: Instant \\ - ageInDays: Double
  \nodepart{third}
};

\node[cls,below=2.4cm of host] (vuln) {
  \textbf{VulnInfo}
  \nodepart{second} - id: Long \\ - qidValue: String \\ - severity: Severity \\ - qds: Integer \\ - firstFound: Instant \\ - lastfixed: Instant \\ - status: String
  \nodepart{third}
};

\node[cls,right=3.6cm of host] (cmdb) {
  \textbf{CmdbCi}
  \nodepart{second} - ciId: String \{PK\} \\ - ip: String \\ - role: Role \\ - env: Env \\ - cmdbOwner: String \\ - soc: boolean \\ - antivirus: boolean \\ - standardPatching: boolean \\ - exception: boolean
  \nodepart{third}
};

\node[cls,below=2.4cm of cmdb] (config) {
  \textbf{RiskConfig}
  \nodepart{second} - wCritical: Double \\ - wHigh: Double \\ - wMedium: Double \\ - wLow: Double \\ - iExternal: Double \\ - mSoc: Double \\ - mAv: Double \\ - mPatch: Double \\ - blackMinArs: Double \\ - redMinArs: Double \\ - orangeMinArs: Double \\ - yellowMinArs: Double \\ - slaBlackDays: Integer \\ - slaRedDays: Integer \\ - slaOrangeDays: Integer \\ - slaYellowDays: Integer \\ - slaGreenDays: Integer
  \nodepart{third}
};

\node[cls,right=2.6cm of cmdb] (wconf) {
  \textbf{RiskWeightConfig}
  \nodepart{second} - category: WeightCategory \\ - key: String \\ - weight: Double
  \nodepart{third}
};

\node[cls,below=2.4cm of wconf,text width=3.4cm] (snap) {
  \textbf{RiskSnapshot}
  \nodepart{second} - computedAt: Instant \\ - blackCount: Long \\ - redCount: Long \\ - orangeCount: Long \\ - yellowCount: Long \\ - greenCount: Long \\ - avgArs: Double
  \nodepart{third}
};

% Association (composition Host -> VulnInfo)
\draw[rel] (host) -- node[left,font=\scriptsize]{1} node[below left=0.05cm,font=\scriptsize]{0..*} (vuln)
  node[midway,right=0.1cm,font=\scriptsize]{vulns};

% Dependances (fleche pointillee vers la classe dont on depend)
\draw[dep] (host) -- node[above,align=center,font=\scriptsize]{reconciliation\\ par ip} (cmdb);
\draw[dep] (host) to[bend right=10] node[right=0.1cm,font=\scriptsize]{parametre} (config);
\draw[dep] (cmdb) to[bend left=15] node[above,font=\scriptsize]{parametre} (wconf);
\draw[dep] (snap.west) to[bend right=15] node[sloped,above,font=\scriptsize,pos=0.4]{instantane de} (host.south);

\node[font=\scriptsize\itshape,dxcgray,below=0.15cm of vuln,text width=4.4cm] {
  Enumerations associees : Severity (CRITICAL, HIGH, MEDIUM, LOW), Role, Env, Os
};
\node[font=\scriptsize\itshape,dxcgray,below=0.15cm of snap,text width=4.2cm] {
  RiskClass (BLACK..GREEN), SlaStatus (ON\_TRACK, VIOLATED, ACCEPTED\_RISK)
};

% Notes UML (rectangle a coin plie), reliees par une ligne pointillee sans fleche
\node[note,above=0.9cm of host] (hostnote) {Attributs calcules par \texttt{RiskCalculationService}, non persistes tels quels.};
\draw[fill=dxcgray!25,draw=dxcgray] (hostnote.north east) -- ++(-0.25,0) -- ++(0.25,-0.25) -- cycle;
\draw[dxcgray] ($(hostnote.north east)+(-0.25,0)$) -- ($(hostnote.north east)+(0,-0.25)$);
\draw[noteline] (hostnote.south) -- (host.north);

\node[note,below=1.4cm of config] (confnote) {Table de configuration a ligne active unique (singleton applicatif).};
\draw[fill=dxcgray!25,draw=dxcgray] (confnote.north east) -- ++(-0.25,0) -- ++(0.25,-0.25) -- cycle;
\draw[dxcgray] ($(confnote.north east)+(-0.25,0)$) -- ($(confnote.north east)+(0,-0.25)$);
\draw[noteline] (confnote.north) -- (config.south);

\end{tikzpicture}
\end{adjustbox}
\caption{Diagramme de classes simplifie du domaine}
\label{fig:class-diagram}
\end{figure}

\subsection{Le choix de la reconciliation par cle applicative}
\label{sec:reconciliation}

Qualys et la CMDB sont deux sources gerees par des equipes et des rythmes distincts :
un hote peut avoir ete scanne sans encore figurer dans la CMDB (ou inversement).
Modeliser une cle etrangere stricte entre \texttt{Host} et \texttt{CmdbCi} aurait
oblige a rejeter ou differer l'ingestion de l'une des deux sources en cas
d'incoherence temporaire. Le choix retenu -- reconcilier les deux entites par
adresse IP au moment du calcul (\texttt{Map<String, CmdbCi>} indexee par IP dans
\texttt{RiskCalculationService}) -- decouple totalement les deux pipelines
d'ingestion et degrade proprement : un hote sans CI correspondant n'est simplement
pas score (il est compte comme \emph{Shadow IT}, cf. \cref{chap:visualisation}).

\section{Diagrammes de sequence}

\subsection{Ingestion d'un export et recalcul du risque}

Le \cref{fig:seq-ingestion} decrit le scenario le plus important du systeme : de
l'upload d'un fichier a la mise a jour du tableau de bord chez tous les clients
connectes. Deux caracteristiques architecturales sont visibles sur ce diagramme :

\begin{itemize}
  \item \textbf{L'asynchronisme a deux niveaux} : le stockage du fichier
  (\texttt{@Async}) et le lancement du job (\texttt{@Async}) rendent la main
  immediatement a l'appelant HTTP, evitant de bloquer un thread de requete pendant
  un traitement potentiellement long.
  \item \textbf{Le decouplage par evenements apres commit} : le recalcul du risque et
  la diffusion SSE ne sont pas appeles directement -- ils reagissent a des
  evenements (\texttt{IngestionCompletedEvent},
  \texttt{RiskRecalculationCompletedEvent}) publies et consommes en phase
  \texttt{AFTER\_COMMIT}, garantissant que le listener ne lit jamais un etat de base
  de donnees non encore valide par la transaction qui l'a produit.
\end{itemize}

\begin{figure}[H]
\centering
\begin{adjustbox}{width=\textwidth}
\begin{tikzpicture}[
  part/.style={draw,minimum width=2.5cm,minimum height=0.8cm,align=center,font=\scriptsize,fill=dxcblue!10},
  lifeline/.style={dashed,dxcgray},
  act/.style={fill=dxcblue!25,draw},
  msg/.style={-{Latex[length=2mm]},font=\scriptsize},
  amsg/.style={-{Latex[length=2mm]},font=\scriptsize,dashed},
  ret/.style={-{Latex[length=1.6mm]},dashed,font=\scriptsize}
]
\def\yA{0}
\def\ytop{0}
\def\ybot{-15.6}

% Participants (x positions)
\node[part] (A) at (0,\ytop)   {:Administrateur};
\node[part] (B) at (3.0,\ytop) {UploadController};
\node[part] (C) at (6.0,\ytop) {TempFileStorage\\Service};
\node[part] (D) at (9.0,\ytop) {JobService\\+ Spring Batch};
\node[part] (E) at (12.0,\ytop){RiskCalculation\\Service};
\node[part] (F) at (15.0,\ytop){DashboardSse\\Controller};
\node[part] (G) at (18.0,\ytop){:Navigateur};

\foreach \n in {A,B,C,D,E,F,G}{
  \draw[lifeline] (\n.south) -- ++(0,\ybot+1.2);
}

% Messages
\draw[msg] (0,-1) -- (3,-1) node[midway,above,font=\scriptsize]{1: POST /upload/\{type\} (fichier)};
\draw[amsg] (3,-1.6) -- (6,-1.6) node[midway,above]{2: store(file, type) [\texttt{@Async}]};
\node[font=\scriptsize] at (6,-2.2) {3: detecterMimeType() ; sha256(); };
\draw[ret] (6,-2.8) -- (3,-2.8) node[midway,above]{4: id fichier stocke};

\draw[msg] (0,-3.6) -- (3,-3.6) node[midway,above]{5: POST /joblaunch};
\draw[amsg] (3,-4.2) -- (9,-4.2) node[midway,above]{6: launchJob(type) [\texttt{@Async}]};
\node[font=\scriptsize] at (9,-4.9) {7: wipeForReingest(type)};
\node[font=\scriptsize,align=center,text width=6.4cm] at (9,-5.7) {8: JobOperator.start(job) $\to$ chunks de 100\\ (StaxEventItemReader $\to$ Processor $\to$ JpaItemWriter)};
\node[font=\scriptsize,dxcblue] at (9,-6.7) {\emph{publie IngestionCompletedEvent (AFTER\_COMMIT)}};
\draw[ret] (9,-7.2) -- (3,-7.2) node[midway,above]{9: BatchStatus.COMPLETED};

\draw[amsg] (9,-8.0) -- (12,-8.0) node[midway,above]{10: onIngestionCompleted() $\to$ recomputeAll()};
\node[font=\scriptsize,align=center,text width=5.6cm] at (12,-9.0) {11: pour chaque Host (parallele) : \\ computeAcs() \textbar\textbar\ computeWeightedQds() \\ $\to$ ARS, classe, statut SLA};
\node[font=\scriptsize] at (12,-10.1) {12: enregistrer RiskSnapshot};
\node[font=\scriptsize,dxcblue] at (12,-10.8) {\emph{publie RiskRecalculationCompletedEvent}};

\draw[amsg] (12,-11.4) -- (15,-11.4) node[midway,above]{13: onRecalculationCompleted()};
\node[font=\scriptsize] at (15,-12.1) {14: computeMetrics()};
\draw[msg] (15,-12.8) -- (18,-12.8) node[midway,above]{15: SSE event "metrics" (JSON)};
\node[font=\scriptsize,align=center,text width=3cm] at (18,-13.6) {16: Chart.js redessine les graphiques};

\end{tikzpicture}
\end{adjustbox}
\caption{Diagramme de sequence -- ingestion d'un export et recalcul du risque}
\label{fig:seq-ingestion}
\end{figure}

\subsection{Consultation de l'assistant conversationnel}

Le second flux notable est celui de l'assistant RAG. Contrairement a une architecture
RAG classique qui recupere des passages de texte par similarite vectorielle avant
generation, l'agent implemente ici s'appuie sur du \emph{tool calling} : le modele
decide lui-meme, en fonction de la question, quel outil interroger (par exemple
« lister les actifs en classe Black »), l'outil execute une requete reelle sur la
base de donnees, et son resultat est reinjecte dans le contexte avant generation de
la reponse finale, diffusee en flux (\texttt{Flux<String>}) au fur et a mesure de sa
production par le modele local.

\begin{figure}[H]
\centering
\begin{adjustbox}{width=0.92\textwidth}
\begin{tikzpicture}[
  part/.style={draw,minimum width=2.6cm,minimum height=0.8cm,align=center,font=\scriptsize,fill=dxcblue!10},
  lifeline/.style={dashed,dxcgray},
  msg/.style={-{Latex[length=2mm]},font=\scriptsize},
]
\node[part] (U) at (0,0)   {:Analyste};
\node[part] (CTL) at (3.2,0) {RagChatController};
\node[part] (AG) at (6.6,0) {AiAgentService\\ (Langchain4j)};
\node[part] (TL) at (10.0,0) {AiAgentTool};
\node[part] (LLM) at (13.2,0) {Ollama\\ (mistral:7b)};

\foreach \n in {U,CTL,AG,TL,LLM}{
  \draw[lifeline] (\n.south) -- ++(0,-5.6);
}

\draw[msg] (0,-1) -- (3.2,-1) node[midway,above]{1: chat(question)};
\draw[msg] (3.2,-1.6) -- (6.6,-1.6) node[midway,above]{2: chat(question)};
\draw[msg] (6.6,-2.4) -- (10.0,-2.4) node[midway,above]{3: appel outil (si necessaire)};
\draw[msg] (10.0,-3.0) -- (6.6,-3.0) node[midway,above]{4: resultat (donnees BD)};
\draw[msg] (6.6,-3.7) -- (13.2,-3.7) node[midway,above]{5: prompt + contexte outil};
\draw[msg] (13.2,-4.4) -- (0,-4.4) node[midway,above]{6: reponse en flux (streaming)};

\end{tikzpicture}
\end{adjustbox}
\caption{Diagramme de sequence -- interrogation de l'assistant RAG}
\label{fig:seq-rag}
\end{figure}

\section{Choix technologiques}

\begin{longtable}{@{}p{3.4cm}p{9.8cm}@{}}
\toprule
\textbf{Composant} & \textbf{Technologie retenue} \\
\midrule
\endhead
Langage / plateforme & Java 17, Spring Boot \\
Persistance & Spring Data JPA (Hibernate), PostgreSQL 15 \\
Ingestion des fichiers & Spring Batch (traitement par blocs de 100 elements),
lecture XML par flux via StAX (\texttt{StaxEventItemReader}) et marshalling
JAXB (\texttt{Jaxb2Marshaller}, Spring OXM) \\
Detection de type de fichier & Apache Tika (verification du type MIME reel,
independante de l'extension declaree) \\
Mapping objet & MapStruct (DTO $\leftrightarrow$ entites) \\
Asynchronisme & \texttt{@Async} sur des \texttt{Executor} dedies (E/S vs. calcul),
\texttt{CompletableFuture}, bus d'evenements Spring
(\texttt{ApplicationEventPublisher}, \texttt{@TransactionalEventListener}) \\
Interface utilisateur & Rendu cote serveur Thymeleaf, interactions dynamiques HTMX
(en-tete \texttt{HX-Request} pour ne renvoyer que le fragment de page pertinent) \\
Tableau de bord & Server-Sent Events (\texttt{SseEmitter}) cote serveur, Chart.js
cote client (avec le plugin de diagramme de Sankey) \\
Assistant conversationnel & Langchain4j (\texttt{@AiService}, tool calling),
Ollama comme moteur d'inference local (modele \texttt{mistral:7b}) \\
Conteneurisation & Docker, Docker Compose (backend, base de donnees, moteur LLM) \\
\bottomrule
\end{longtable}
